\documentclass[oneside,final,14pt]{extreport}
\usepackage[english,russian]{babel}
\usepackage[T2A]{fontenc}
\usepackage{cmap}
\usepackage{graphicx}
\usepackage{amsmath}
\usepackage{amssymb}
\usepackage[a4paper, left=1cm, right=1cm, top=2cm, bottom=2cm]{geometry}
\usepackage{multirow}
\graphicspath{ {images/} }

\title{\textbf{Вопросы ЧГК в КиМ ноябрь 2026}}
\author{Тараканов Павел, Дитятьева Мария}

\begin{document}

\maketitle

\section*{Вопросы ЧГК в КиМ 2026-2027}
\begin{enumerate}
    \item ОН — один из величайших советских маршалов. Однако первые свои награды ОН получил еще в первую мировую. Сослуживцы шутили над молодым ефрейтором: «ОН ЕГО получил», намекая на самую популярную награду империи. Назовите ЕГО фамилию

    Ответ: Жуков

    Комментарий: Георгий Константинович Жуков участвовал в Первой Мировой войне, где получил два Георгиевских креста, в народе называемых просто Георгиями.

    \item В 1701 году, Пётр 1, потеряв в Нарвском сражении всю артиллерию, приказал восполнить эти потери. Согласно распространённой легенде, он издал указ, повелевавший изымать у монастырей определённое количество ИХ. На самом деле ОНИ не годились для изготовления пушек, потому что имели немного другой состав. В указе имелись в виду не ОНИ, а запасы металла, из которых ОНИ изготавливались. Что такое "ОНИ"?

    Ответ: колокола

    Комментарий: Пушки, как и колокола, отливали из бронзы, однако в пушечной бронзе было меньше олова, чем в колокольной. В то время монастыри были не только местами, где проживали и молились монахи, а ещё и крепостями, в которых имелись пушки, чаще всего применявшиеся при подавлении крестьянских мятежей. Кроме непосредственно пушек, в монастырях были запасы так называемой "колокольной меди", которая использовалась для литья как колоколов, так и пушек.

    \item ОНИ— легкая пехота  18-го и 19-го веков. В годы линейной тактики лишь ОНИ  действовали в рассыпном строю и были единственным родом войск, подходящим для горной местности Кавказа. Действуя малыми группами, они громили огромные отряды персов. Например, 17-й ТАКОЙ полк, состоявший из 493 человек, две недели держался сорока пяти тысяч персов и, потеряв две трети личного состава, не сдался и смог послать за помощью. Сегодня ОН - это профессия человека, который отвечает за охрану лесной фауны. Назовите ИХ словом, пришедшим к нам из немецкого языка.

    Ответ: егерь/егеря

    Комментарий. Ответ - егеря, от немецкого Jäger — охотник. 17 егерский полк под  командованием полковника Корягина и майора Котляревского две недели сражался с  персами, почти без еды и воды, в ужасную жару и засуху.

    \item Осенью 1942 года во время боев на подступах к Сталинграду 22-я немецкая танковая дивизия оказалась в критическом положении. Командование планировало нанести контрудар, однако большая часть танков просто не смогла завестись. Виной тому стали маленькие и незаметные «диверсанты». Кто сорвал боевую операцию вермахта?


    Ответ: Мыши.

    Комментарий: Танки дивизии долгое время стояли в укрытиях, замаскированные соломой. Полевые мыши забрались внутрь боевых машин в поисках тепла и провианта и полностью перегрызли электропроводку, выведя технику из строя без единого выстрела.

    \item В 1939 году во время боев на реке Халхин-Гол против японских войск советские танкисты столкнулись с проблемой: бутылки с зажигательной смесью, брошенные японскими пехотинцами, намертво сцеплялись с жалюзи моторных отсеков танков БТ-7, вызывая пожары. Инженеры нашли копеечное решение, установив на жалюзи мелкоячеистую сетку. Однако на передовой сеток не хватало, и танкисты использовали полужидкий продукт из сухпайка, густо обмазывая им решетки. Какой продукт выполнял роль тушителя?

    Ответ: Сгущенное молоко (сгущенка).

    Комментарий: Густая и липкая сгущенка при нагревании от двигателя превращалась в плотную обуглившуюся корку, которая не давала горящей горючей смеси затекать внутрь моторного отсека, выигрывая время для экипажа.

    \item В 1916 году генерал Алексей Брусинов совершил знаменитый прорыв австро-венгерского фронта. Новаторство заключалось в том, что вместо одного мощного удара (который враг легко вычислял и подтягивал резервы) русская армия нанесла удар сразу в нескольких десятках мест одновременно. В штабах противника началась паника, так как они не могли понять, где именно находится... Что?

    Ответ: Главный удар.

    Комментарий: Из-за одновременного наступления на всем протяжении Юго-Западного фронта австрийское командование не смогло определить направление главного удара и распылило свои резервы, что и привело к краху их обороны.

    \item У Вильгельма 2, кайзера Германской Империи, в годы ПМв было элитное военное подразделение — стальная дивизия. Однажды немцы столкнулись с подобными себе - русский ТАКОЙ дивизией, под командованием ЕГО. Над немецкими позициями появился плакат, на котором было написано: «Ваше русское ОНО не хуже нашей германской стали, а всё же мы вас разобьём». Русские ответили: «А ну, попробуй!» За 4 дня немцы предприняли 42 атаки на русские позиции и, не добившись успеха, отступили с огромными потерями. А командир русской дивизии через много лет, в ходе Великой Отечественной войны, на свои средства направил в СССР вагон с медикаментами. Вагон приняли, скрыв имя отправителя. Назовите, КАКОЙ была русская дивизия, или фамилию русского командира.

    \item Художник и дизайнер ЙЕрун ван дер Мост разработал программу, которая с помощью вайб-кОдинга меняет цвета ячеек Excel в реальном времени. Недавно в Таллинском музее в рамках экспозиции, посвящённой Excel, был показан фильм. Назовите этот фильм.

    Ответ: «Матрица». Зачёт: по упоминанию «Матрицы».

    Комментарий. Видим код – говорим «Матрица».

    \item ОНИ и великолепная память помогли Николаю Пржевальскому добыть средства на первую экспедицию. Некоторые другие ОНИ после путешествий Пржевальского изменились, причём в лучшую сторону. Назовите ИХ одним словом.

    Ответ: карты.

    Комментарий. Пржевальский был успешным картёжником, отлично запоминал
    комбинации. Добыв игрой необходимую сумму, карты навсегда оставил. Путешествия Пржевальского позволили составить либо уточнить географические карты Восточной и Центральной Азии.

    \item Пакистанский ОН известен меньше, чем ОН соседней страны, и гораздо меньше, чем ОН американский. Поскольку эта индустрия наиболее развита в Лахоре, пакистанскому ЕМУ, по аналогии с двумя другими, дали неофициальное прозвище Лол [ЛОЛ]… Напишите полностью слово, сокращённое в предыдущем предложении.

    Ответ: Лолливуд.

    Комментарий. ОН — кинематограф. По аналогии с более известными американским Голливудом и индийским Болливудом (названным по городу Бомбею), кинематограф пакистанского Лахора называют Лолливудом. Кстати, на заре пакистанского кинематографа женские роли играли мужчины.

    \item Развивая известное сравнение Земли с НЕЮ, РОберт ЗУбрин заметил, что тогда Солнечная система – это школьный двор. ОНА входит в название изобретения, статья Википедии о котором сопровождена несколькими gif-изображениями . Назовите ЕЁ одним словом.

    Ответ: колыбель.

    Комментарий. РОберт ЗУбрин как-то заметил: «Если Земля – колыбель человечества, то Солнечная система – школьный двор». В статье Википедии «Колыбель НьЮтона» для наглядности показано, как работает система шаров.


    \item Математический анекдот. Математика спрашивают: "Есть ли крылья у слона?" - "Есть, - отвечает математик, но они..." Закончите фразу двумя словами!

    Ответ: "... равны нулю".

    Комментарий: комментарии излишни

    \item Станислав Ежи Лец спрашивает: «Что ты скажешь на это, физика? В отношениях между людьми трения ведут к … . А к чему ведут трения?

    Ответ: К охлаждению

    Комментарий: в физике при трении происходит выделение тепла и нагревание тел.

    \item В 1912 году в газетах появилась сенсационная фотография: на палубе корабля стоят хорошо одетые господа, а над ними в воздухе завис странный аппарат. Это был один из первых авиационных трюков для прессы. Однако физик, взглянув на фото, сразу сказал бы, что это монтаж. Какая физическая несообразность выдала подделку?

    Ответ: Отсутствие ветра / Прически и одежда.

    Комментарий: Если бы дирижабль или самолет висел прямо над палубой на малой высоте, винты создавали бы мощный поток воздуха, который раздувал бы одежду и волосы людей. На фото люди стояли со спокойными прическами, значит, снимок смонтирован.

    \item Голландский физик Гендрик Лоренц нашел соотношение между плотностью вещества и его преломляющей способностью. Несколько ранее эту же формулу, но совершенно другим путем получил датский физик. За формулой утвердилось двойное название. Открытие этой формулы является любопытным случаем с точки зрения теории вероятности. Как называется эта формула?

    Ответ: Формула Лоренца-Лоренца

    Комментарий: фамилия датского физика тоже была Лоренц

    \item Это теория, в названии которой есть 2 слова: существительное и прилагательное. Одно из них связано с известным героем русских народных сказок, а другое - с самыми мощными вычислительными машинами на Земле.

    \item В 2021 году человек под ником IlluminatiPirate опубликовал конспирологическую теорию, в которой утверждал, что, начиная с 2016—2017 годов, основная часть контента в интернете создаётся ботами и системами ИИ ,что «живой» интернет почти исчез. Исследователи считают ее метафорой, полагая, что люди просто боятся неконтролируемого изменения интернет-пространства. Какое название носит эта теория?

    Ответ: теория мертвого интернета

    \item Вопрос: Космодесантники из фантастической вселенной Warhammer 40000 должны выдерживать огромные нагрузки, поэтому в начале подготовки у космодесантника появляется ОНО. ИМ также называют икроножные мышцы. Назовите ЕГО двумя словами, одно из которых - числительное.

    Ответ: Второе сердце.

    Комментарий: У космодесантников имплантируется дополнительный орган — гемостата (второе сердце), помогающий перекачивать кровь при тяжелых ранениях и нагрузках.

    \item В традиционной корейской медицине было принято добавлять солодку во все лекарства. Корейское выражение “как сироп солодки в аптеке” примерно соответствует нашему фразеологизму, в котором фигурируют АЛЬФА и БЕТА. Назовите АЛЬФУ и БЕТУ словами, которые заканчиваются одинаково. А если вы не полностью поняли вопрос – сдайте хотя бы БЕТУ.

    Ответ: бочка, затычка.

    Зачет: затычка.

    Комментарий: корейцы так говорят про человека, без которого не обходится ни одно дело. Это примерный аналог нашего выражения “в каждой бочке затычка”. Версию, которая сдаётся из общих соображений, без полного понимания сути вопроса, называют “затычкой”.
\end{enumerate}

\end{document}
